% The resume of Roy T. Smart, built on the "Medium Length Professional CV"
% template by Trey Hunner (http://www.treyhunner.com/), from
% http://www.LaTeXTemplates.com. The layout lives in resume.cls.

\documentclass{resume}

\usepackage[left=0.5in,top=0.4in,right=0.5in,bottom=0.4in]{geometry}
\usepackage{fontawesome5}
\usepackage{enumitem}
\usepackage{multicol}
\usepackage{tabularx}
\usepackage{xcolor}
% biblatex-apa lists up to 20 authors, so long author lists are truncated with
% `and others` in sources.bib.
\usepackage[style=apa, backend=biber]{biblatex}
\usepackage[colorlinks, urlcolor=blue!50!black, linkcolor=black]{hyperref}

\addbibresource{sources.bib}

% Link the draft of an unpublished article by name rather than by its full url.
\DeclareFieldFormat[unpublished]{url}{\href{#1}{Draft}}

% Bold my name in author lists, using the `highlight` annotation in sources.bib.
\renewcommand*{\mkbibnamegiven}[1]{\ifitemannotation{highlight}{\textbf{#1}}{#1}}
\renewcommand*{\mkbibnamefamily}[1]{\ifitemannotation{highlight}{\textbf{#1}}{#1}}

% Sub-bullets beneath a bullet of an `rSubsection`.
\newlist{details}{itemize}{1}
\setlist[details]{topsep=0pt, noitemsep, leftmargin=*, label={--}}

% A list of citations or talks, set ragged right so that author lists are not
% stretched to fill the line.
\newenvironment{entries}{%
	\small\raggedright
	\begin{itemize}[leftmargin=*, topsep=2pt, itemsep=1pt, parsep=0pt]%
}{%
	\end{itemize}%
}

% A table with bold labels on the left, which fills the width of the line.
\newenvironment{labeled}{%
	\small
	\tabularx{\linewidth}{@{} >{\bfseries}l @{\hspace{3ex}} X @{}}%
}{%
	\endtabularx
}

\setlength{\multicolsep}{0pt}

% An entry with a title on the left and dates on the right.
\newcommand{\dated}[2]{\textbf{#1} \hfill {\fontseries{b}\selectfont #2}}

\name{Roy T. Smart}

\address{Department of Physics, Montana State University, Bozeman, MT}
\address{%
	\faEnvelope[regular]~\href{mailto:roytsmart@gmail.com}{roytsmart@gmail.com} \\
	\faGlobe~\href{https://roytsmart.github.io}{roytsmart.github.io} \\
	\faGithub~\href{https://github.com/roytsmart}{roytsmart} \\
	\faOrcid~\href{https://orcid.org/0000-0002-9997-5515}{0000-0002-9997-5515}%
}

\begin{document}

\begin{rSection}{Research Interests}
	\begin{multicols}{2}
		\begin{itemize}[leftmargin=*, noitemsep, topsep=0pt]
			\item Solar transition region explosive events and jets
			\item Snapshot imaging spectroscopy
			\item EUV instrument design and optical modeling
			\item Ultraviolet response of silicon sensors
			\item Inverse problems and machine learning
			\item Open-source scientific software
		\end{itemize}
	\end{multicols}
\end{rSection}

\begin{rSection}{Education}

	\dated{Ph.D. in Physics, Montana State University}{Aug 2015 -- May 2027 (expected)}
	\\
	Advisor: Charles Kankelborg

	\dated{B.S. in Physics, Minor in Computer Science, Montana State University}{Aug 2011 -- May 2015}

\end{rSection}

\begin{rSection}{Research Experience}

	\begin{rSubsection}{Department of Physics, Montana State University}{Jan 2016 -- Present}{Graduate Research Assistant, advised by Charles Kankelborg}{Bozeman, MT}
		\small
		\item \textbf{EUV Snapshot Imaging Spectrograph (ESIS).}
		A NASA sounding rocket computed tomography imaging spectrograph, which measures Doppler shifts across an 11.5 arcminute field of the transition region in a single exposure.
		\begin{details}
			\item \textbf{Instrument paper.} Leading the instrument paper, a reproducible article generated from the instrument model.
			\item \textbf{Optical model.} Developed the optical model of ESIS in Python, which predicts its distortion, vignetting, point-spread function, and effective area from measurements of its coatings, filters, and gratings.
			\item \textbf{Flight preparation.} Developed the analysis used to align and focus the optics, generated the baffle apertures from the optical model, and measured the wavefront error and roughness of the primary mirror and gratings with phase-shifting interferometry.
			\item \textbf{Calibration.} Built the pipeline that calibrates the flight images, including the characterization of the gain, dark current, and linearity of the CCD cameras built by NASA Marshall Space Flight Center.
			\item \textbf{Science.} Inverted the flight images into spectral cubes using convolutional neural networks and the multiplicative algebraic reconstruction technique. Leading a study of a transition region jet which combines ESIS Doppler shifts with the magnetic topology from SDO/HMI and differential emission measures from SDO/AIA and Hinode/XRT.
			\item \textbf{ESIS-II.} Helped design the optics of the second flight (2027), which adds channels and new spectral lines.
		\end{details}
		\item \textbf{Silicon sensors in the ultraviolet.}
		Proposed partial charge collection as the reason that Hubble/WFC3 and IRIS measure less noise in the ultraviolet than photon shot noise predicts, with a model that agrees with both.
		Measured the depth-dependent charge-diffusion kernel of all four IRIS CCDs in orbit, using energetic particles that cross the sensors at glancing angles.
		\item \textbf{Full-sun Ultraviolet Rocket Spectrometer (FURST).}
		Contributed to the optical design and tolerance analysis of this NASA sounding rocket spectrograph of the Sun as a star.
		Modeled the instrument as built, from measurements of its flight grating, coatings, and filter, to predict its effective area and spectral resolution for the instrument paper.
		\item \textbf{Multi-Order Solar EUV Spectrograph (MOSES).}
		Measured Doppler shifts in MOSES-II images by tracking features between spectral orders, and trained neural networks on IRIS observations to invert MOSES images.
		\item \textbf{IRIS planner.} Chose the targets and observing programs of the Interface Region Imaging Spectrograph, and coordinated its observations with other observatories.
	\end{rSubsection}

	\begin{rSubsection}{Department of Physics, Montana State University}{Sep 2012 -- Aug 2015}{Undergraduate Research Assistant, advised by Charles Kankelborg}{Bozeman, MT}
		\small
		\item \textbf{MOSES-II.} Wrote the flight software that commanded MOSES-II and handled its data, and analyzed the thermal behavior of an optical mount.
	\end{rSubsection}

\end{rSection}

\begin{rSection}{Sounding Rocket Campaigns}
	\textbf{ESIS}, White Sands Missile Range, NM: integration, testing, and flight operations \hfill {\fontseries{b}\selectfont Sep 30, 2019}
	\\
	\textbf{MOSES-II}, White Sands Missile Range, NM: flight software and range operations \hfill {\fontseries{b}\selectfont Aug 27, 2015}
\end{rSection}

\begin{rSection}{Fellowships}
	\dated{NASA Earth and Space Science Fellowship (NESSF)}{Sep 2017 -- Sep 2020}
	\\
	\textit{Neural Networks for Computed Tomography Imaging Spectroscopy of the Solar Atmosphere}
\end{rSection}

\begin{rSection}{Publications}
	\textbf{Published and in press}
	\begin{entries}
		\item \fullcite{LongcopeRidge}
		\item \fullcite{Parker2022}
		\item \fullcite{Johnson2018}
	\end{entries}

	\textbf{In preparation}
	\begin{entries}
		\item \fullcite{SmartESIS}
		\item \fullcite{SmartJet}
		\item \fullcite{SmartNoise}
		\item \fullcite{SmartDiffusion}
	\end{entries}
\end{rSection}

\begin{rSection}{Selected Talks}
	\begin{entries}
		\item The Physical Structure and Evolution of Transition Region Explosive Events Observed with ESIS. \textit{AAS SPD}, 2026.
		\item The Noise Measured by Back-Illuminated Silicon Sensors Modeled Using Partial Charge Collection. \textit{OpTeC} and \textit{SPD}, 2025.
		\item Inverting ESIS Observations of the Solar Transition Region using Convolutional Neural Networks. \textit{IRIS-Hinode}, 2024.
		\item Diagnosing Coronal Heating using Computed Tomography Imaging Spectroscopy. \textit{ISSI team meeting}, 2023.
		\item A Convolutional Neural Network for Inverting Observations from ESIS. \textit{Colloquium, LMSAL}, 2022.
		\item Analyzing ESIS Observations Using Convolutional Neural Networks. \textit{Colloquium, NASA MSFC}, 2020.
		\item Machine Learning Techniques for Computed Tomography Imaging Spectroscopy of the Solar Atmosphere. \textit{AGU}, 2018.
	\end{entries}
	\small
	Also posters at AGU (2019, 2020) and the AAS SPD (2016, 2017). All talks, with slides, are at \href{https://roytsmart.github.io/talks.html}{roytsmart.github.io/talks}.
\end{rSection}

\begin{rSection}{Software}
	\small
	Author and maintainer of open-source Python packages for modeling instruments and analyzing solar data, each released on PyPI with a Zenodo DOI, tested in continuous integration, and documented online. Most are at \href{https://github.com/sun-data}{github.com/sun-data}.

	\begin{labeled}
		\href{https://named-arrays.readthedocs.io/en/stable/}{named-arrays} & Arrays with named axes and automatic propagation of uncertainties \\
		\href{https://optika.readthedocs.io/en/stable/}{optika} & Raytracing and modeling of optical systems, similar to Zemax \\
		\href{https://ctis.readthedocs.io/en/stable/}{ctis} & Inversion of images captured by computed tomography imaging spectrographs \\
		\href{https://esis-mission.github.io/esis/}{esis}, \href{https://furst-optics.readthedocs.io/en/stable/}{furst}, \href{https://msfc-ccd.readthedocs.io/en/stable/}{msfc-ccd} & Models of the ESIS and FURST instruments and the MSFC CCD cameras \\
		\href{https://interface-region-imaging-spectrograph.readthedocs.io/en/stable/}{iris}, \href{https://sdo.readthedocs.io/en/stable/}{sdo}, \href{https://hinode.readthedocs.io/en/stable/}{hinode}, \href{https://utu.readthedocs.io/en/stable/}{utu} & Analysis of IRIS, SDO, and Hinode data, and spectral lines from CHIANTI \\
		\href{https://aastex.readthedocs.io/en/stable/}{aastex} & Journal articles written as Python programs, so every number in the text is computed \\
	\end{labeled}
\end{rSection}

\begin{rSection}{Skills}
	\begin{labeled}
		Optics & Zemax OpticStudio, phase-shifting interferometry, alignment and focus, EUV coatings and filters \\
		Detectors & CCD characterization: gain, photon transfer, dark current, linearity, charge diffusion \\
		Solar data & SDO/AIA and HMI, IRIS, Hinode/XRT, CHIANTI, differential emission measures \\
		Languages & Python, C/C++, CUDA, IDL, \LaTeX \\
		Python & NumPy, SciPy, Astropy, SunPy, Numba, Keras, Matplotlib, pytest, Sphinx, GitHub Actions \\
	\end{labeled}
\end{rSection}

\begin{rSection}{Teaching and Mentoring}
	\textbf{Mentor}, ESIS undergraduate optics team
	\\
	\textbf{Graduate Teaching Assistant}, Physics I laboratory \hfill {\fontseries{b}\selectfont Aug 2015 -- Dec 2015}
\end{rSection}

\begin{rSection}{Personal Interests}
	Alpine skiing, whitewater rafting, and amateur astronomy. \hfill \textit{References are available on request.}
\end{rSection}

\end{document}
